\section{Comparison with Existing Formats}
\label{sec:comparison}

\tesserahash{} occupies a specific niche: a \emph{content identifier} that carries provenance metadata alongside a hash strong enough to outlive deployments. This section places it against adjacent constructs.

\subsection{Hash functions and digest formats}

\begin{center}
\small
\begin{tabularx}{\textwidth}{@{}lXXX@{}}
\toprule
\textbf{Construct} & \textbf{Output bits} & \textbf{PQ posture} & \textbf{Provenance metadata} \\
\midrule
SHA-256                         & 256  & $2^{128}$ classical / $2^{85}$ BHT collision        & none \\
SHA3-512 \citep{fips202}        & 512  & $2^{256}$ classical / $2^{170}$ BHT collision       & none \\
SHAKE-256 \citep{fips202}       & arbitrary & up to $2^{170}$ BHT at 512-bit output           & none \\
BLAKE3                          & 256 (default) & $2^{128}$ classical / $2^{85}$ BHT collision  & none \\
BLAKE2b-512                     & 512  & $2^{256}$ classical / $2^{170}$ BHT collision       & none \\
Git SHA-1                       & 160  & $2^{80}$ classical / $\sim 2^{53}$ BHT collision    & none (broken in 2017) \\
Git SHA-256 (transition)        & 256  & $2^{128}$ classical / $2^{85}$ BHT collision        & none \\
IPFS CIDv1 (multihash)          & variable & depends on chosen hash                          & multihash codec ID \\
\textbf{\tesserahash{} v1}           & \textbf{512} & \textbf{$2^{256}$ classical / $2^{170}$ BHT collision} & \textbf{TS (48b) + CTR (16b) + VAR (8b)} \\
\bottomrule
\end{tabularx}
\end{center}

TesseraHash is not in the same category as ``a hash function''; it is a wire format that bundles SHAKE-256 with provenance. The right way to think about \tesserahash{} is ``IPFS CIDv1 if CIDv1 had committed to one hash and one provenance shape and skipped the multihash codec table.''

\subsection{Identifier formats with embedded hashes}

\begin{center}
\small
\begin{tabularx}{\textwidth}{@{}lllX@{}}
\toprule
\textbf{Format} & \textbf{Size (bytes)} & \textbf{Hash bits} & \textbf{Notes} \\
\midrule
Git commit hash (SHA-1)         & 20 & 160 & Broken; transition to SHA-256 underway \\
Git commit hash (SHA-256)       & 32 & 256 & Forward-compat; no provenance metadata \\
IPFS CID (default)              & 36 & 256 & SHA-256 + codec + version; provenance via outer protocol \\
Docker image digest (SHA-256)   & 32 & 256 & No provenance \\
Subresource Integrity (SHA-512) & 64 & 512 & Inline data-URI; no provenance \\
ULID \citep{ulid}               & 16 & 0 (random) & Identifier, not content-bound \\
\tesseraid{}                          & 41 & 0 (256-bit random) & Identifier, not content-bound \\
\textbf{\tesserahash{} v1}           & \textbf{73} & \textbf{512} & Content-bound + provenance metadata \\
\bottomrule
\end{tabularx}
\end{center}

The deliberate scope difference: \tesseraid{} answers ``which record is this?'' (random; unguessable); \tesserahash{} answers ``which content is this?'' (deterministic in content; collision-resistant). When both questions matter, use both --- the canonical \tesseraid{}-as-AD pattern in \tesseraseal{} (\cref{sec:monorepo}) lets a \tesseraseal{} envelope bind one TesseraID record to one TesseraHash content commitment, and the wire formats stack cleanly.

\subsection{Where \tesserahash{} wins}

\begin{itemize}
\item \textbf{Multi-decade collision margin.} 512-bit SHAKE-256 is the highest commonly-deployed standardized hash. $2^{256}$ classical collision is two orders of magnitude beyond what SHA-256-based formats (Git, Docker, IPFS-default) provide.
\item \textbf{Provenance built in.} The 9-byte header carries enough metadata to make a \tesserahash{} sortable by generation time, traceable to a single generator's counter sequence, and decode-rejectable from cross-project siblings. Bare hash digests carry none of this.
\item \textbf{One primitive across the monorepo.} The same SHAKE-256 implementation that produces \tesserahash{} tails is reused inside \tesseraseal{}'s HKDF schedule (\cref{sec:monorepo}); two primitives audited, one used everywhere.
\item \textbf{Canonical text encodings.} 98-char Base64URL and 117-char Crockford Base32 with explicit padding-bit guards. Bare hex (147 chars) is the default for SHA-512 digests in most ecosystems; \tesserahash{}'s encodings are 33--38\% shorter and the Crockford form sorts by generation time.
\end{itemize}

\subsection{Where \tesserahash{} does not fit}

\begin{itemize}
\item \textbf{If you need an unkeyed-MAC--like property.} Use HMAC-SHA3-256 or KMAC256. \tesserahash{} provides no producer authentication.
\item \textbf{If you need to substitute for SHA-256 inside an existing protocol.} The protocol expects 32 bytes; \tesserahash{} emits 73. Use \texttt{hash\_only(content)} (64-byte tail only) and drop the provenance.
\item \textbf{If you need multihash agility} (IPFS-style). TesseraHash commits to SHAKE-256; adding a second hash requires a new variant byte under the \tessera{} discipline. TesseraHash will not chase the multihash codec table.
\item \textbf{If you need confidentiality of the content.} Use \tesseraseal{} (\cref{sec:monorepo}); \tesserahash{} reveals the content to anyone who can guess it.
\end{itemize}
